\chapter{Implementation}

SEDS is being developed using Agile sprints, with completed sprints delivering authentication, station and zone management, and the interactive map view. The emergency dispatch core (SOS handling, spatial triage, routing, and real-time telemetry) remains planned. The system is approximately 35\% complete.

\section{Preliminary Research}

Prior to implementation, research was conducted on spatial databases, PostGIS geometry indexing, and GIS-based emergency dispatch architectures. Real-time WebSocket patterns and OAuth2 integration approaches were also reviewed.

\section{Current Implementation Status}

The system is approximately 35\% complete. Core modules implemented include authentication, station and zone management, and the interactive map view.

\subsection{Frontend}

\textbf{Landing Page:} The landing page renders an interactive MapLibre GL map displaying all registered emergency stations as points and their service zones as polygons. Users can filter by category (POLICE, FIRE, MEDICAL).

\begin{figure}[H]
  \centering
  \includegraphics[width=\textwidth]{landing-page.png}
  \caption{Landing Page: Map View with Stations and Zones}
  \label{fig:landing}
\end{figure}

\textbf{Log In:} The login page authenticates users and sets a JWT in an HTTP-only cookie, then redirects based on role.

\begin{figure}[H]
  \centering
  \includegraphics[width=\textwidth]{login-page.png}
  \caption{Login Page}
  \label{fig:login}
\end{figure}
\FloatBarrier

\newpage
\textbf{Admin Dashboard:} The admin dashboard lists all stations with their category and location, and allows the admin to add new stations.

\begin{figure}[H]
  \centering
  \includegraphics[width=\textwidth]{admin-dashboard.png}
  \caption{Admin Dashboard: Station Management}
  \label{fig:dashboard}
\end{figure}

\textbf{Add Station Form:} The Add Station dialog allows the admin to specify the station name and category, place a point on the MapLibre GL map, and draw or upload a service zone polygon.

\begin{figure}[H]
  \centering
  \includegraphics[width=\textwidth]{add-station-form.png}
  \caption{Add Station Form: Location Pin and Zone Drawing}
  \label{fig:addstation}
\end{figure}
\FloatBarrier

\subsection{Backend}
The backend is built with Node.js and Express, providing a REST API under \texttt{/api/v1/}:
\begin{itemize}
  \item \textbf{Auth routes:} signup, login, logout, email verification, password reset, Google OAuth2.
  \item \textbf{User routes:} profile management, password change.
  \item \textbf{Admin routes:} station and zone creation (atomic PostGIS transaction), GeoJSON endpoints.
\end{itemize}

\subsection{Database}
PostgreSQL 16 with PostGIS 3.4 runs in Docker. Station locations are stored as \texttt{GEOMETRY(Point, 4326)} and zone boundaries as \texttt{GEOMETRY(Polygon, 4326)} with GIST indexes. The station-zone creation uses an atomic Sequelize transaction with \texttt{ST\_GeomFromGeoJSON} and \texttt{ST\_SetSRID}.

\subsection{Authentication}
Secure authentication is implemented using JWT stored in HTTP-only cookies:
\begin{itemize}
  \item \textbf{Role-Based Access:} Admin, Responder, and Citizen roles enforced by Express middleware.
  \item \textbf{Local Auth:} bcrypt password hashing; email verification via Nodemailer.
  \item \textbf{Google OAuth2:} Passport-style Google OAuth2 flow via \texttt{googleapis} SDK.
\end{itemize}

\section{Tools Used}
\begin{itemize}
  \item \textbf{VS Code:} Primary code editor.
  \item \textbf{Node.js 20 + TypeScript:} Backend runtime and language.
  \item \textbf{Next.js 16:} Frontend framework.
  \item \textbf{PostgreSQL 16 + PostGIS 3.4:} Spatial database.
  \item \textbf{Docker + Docker Compose:} Containerization and orchestration.
  \item \textbf{MapLibre GL JS:} WebGL map rendering library.
  \item \textbf{Postman:} API testing.
  \item \textbf{GitHub:} Version control.
\end{itemize}

\section{Testing Overview}

API endpoints were tested using Postman. Manual testing covered form submissions, map interactions, and role-based redirects.

\begin{table}[H]
\centering
\caption{API Test Cases}
\rowcolors{1}{gray!10}{white}
\begin{tabularx}{\textwidth}{|X|m{1.6cm}|X|m{1.2cm}|}
\hline
\textbf{Endpoint} & \textbf{Method} & \textbf{Expected Outcome} & \textbf{Result} \\
\hline
/api/v1/auth/signup & POST & Creates user, returns 201 & Pass \\
\hline
/api/v1/auth/login & POST & Returns JWT cookie, 200 & Pass \\
\hline
/api/v1/auth/logout & POST & Clears cookie, 200 & Pass \\
\hline
/api/v1/auth/verify-email & POST & Activates account, 200 & Pass \\
\hline
/api/v1/auth/forgot-password & POST & Sends reset email, 200 & Pass \\
\hline
/api/v1/user/profile & GET & Returns profile, 200 & Pass \\
\hline
/api/v1/admin/stations & POST & Creates station + zone, 201 & Pass \\
\hline
/api/v1/admin/stations & GET & Lists all stations, 200 & Pass \\
\hline
/api/v1/admin/stations/geojson & GET & Returns GeoJSON, 200 & Pass \\
\hline
/api/v1/admin/zones/geojson & GET & Returns GeoJSON, 200 & Pass \\
\hline
/api/v1/admin/stations & POST (citizen role) & 403 Forbidden & Pass \\
\hline
\end{tabularx}
\end{table}
